\documentclass[11pt]{article}
\usepackage[a4paper,margin=27mm]{geometry}
\usepackage{amsmath,amssymb,amsthm,xcolor,booktabs,array}
\newcommand{\PROVED}{\textcolor{green!40!black}{\textbf{PROVED}}}
\newcommand{\OPEN}{\textcolor{red!70!black}{\textbf{OPEN}}}
\newcommand{\AUDIT}{\textcolor{orange!70!black}{\textbf{AUDIT}}}
\newtheorem{theorem}{Theorem}
\newtheorem{corollary}{Corollary}
\title{The completed inner scattering model and its Tate centering (v115)}
\date{13 August 2026}

\begin{document}
\maketitle

\section{A cross-place inner function}

Let $\xi$ be the completed entire Riemann function, normalized by
$\xi(s)=\xi(1-s)$ and $\xi(\overline s)=\overline{\xi(s)}$.  Put
\begin{equation}
                         \Theta(z)=\frac{\xi(z)}{\xi(z+1)}.    \tag{1}
\end{equation}

\begin{theorem}[Completed meromorphic inner scattering --- \PROVED]
The function $\Theta$ is analytic in the right half-plane, has unimodular
boundary values on $i\mathbb R$, and satisfies
\begin{equation}
                         \Theta(-z)=\Theta(z)^{-1}.             \tag{2}
\end{equation}
Its zeros in the right half-plane are exactly the nontrivial zeros $\rho$
of $\xi$, with their multiplicities.
\end{theorem}

\begin{proof}
The denominator has no zero for $\Re(z+1)>1$.  On $z=i\tau$, the functional
equation and reality give
\[
 \xi(i\tau)=\xi(1-i\tau)=\overline{\xi(1+i\tau)},
\]
so $|\Theta(i\tau)|=1$ away from a null set, and continuation supplies the
boundary values.  Also
\[
 \Theta(-z)=\frac{\xi(-z)}{\xi(1-z)}
            =\frac{\xi(1+z)}{\xi(z)}=\Theta(z)^{-1}.
\]
The zero statement follows directly from (1), since every nontrivial zero
has $0<\Re\rho<1$ and the denominator is zero-free in the domain.  The
Hadamard product pairs every zero with its reflected pole across the
boundary; together with the boundary modulus this is the meromorphic-inner
factorization of (1).
\end{proof}

Thus (1) mixes the Euler and Gamma data before scalar continuation cancels
the finite-place labels.  Let
\begin{equation}
 \mathscr K_\Theta=H^2(\mathbb C_+)\ominus\Theta H^2(\mathbb C_+)          \tag{3}
\end{equation}
be its positive model Hilbert space.

\section{The model semigroup and the spectral characters}

For $t\ge0$, multiplication by $e^{-tz}$ is a contraction on
$H^2(\mathbb C_+)$.  Let $T_t$ be its compression to $\mathscr K_\Theta$.
If $k_\rho$ is the reproducing kernel at a zero $\rho$ of $\Theta$, then
\begin{equation}
                         T_t^*k_\rho=e^{-t\overline\rho}k_\rho.            \tag{4}
\end{equation}

\begin{theorem}[Positive Hilbert realization of finite zero blocks ---
\PROVED]
Every finite collection of zero kernels and their derivative kernels spans
a finite-dimensional $T_t^*$-invariant Hilbert space.  Its characteristic
values are $e^{-t\overline\rho}$, with the zero multiplicities.  Hence its
finite trace is
\begin{equation}
                         \operatorname{Tr}T_t^*
 =\sum_\rho m_\rho e^{-t\overline\rho}                         \tag{5}
\end{equation}
over the selected finite multiset.
\end{theorem}

\begin{proof}
Equation (4) is the reproducing-kernel identity
$\langle f,T_t^*k_\rho\rangle=(T_tf)(\rho)
=e^{-t\rho}f(\rho)$.  Differentiating at a multiple zero gives the usual
upper-triangular derivative-kernel block, whose diagonal repeats the same
character.  Finite-dimensional trace is the sum of diagonal characters.
\end{proof}

This supplies a genuine positive cross-place Hilbert realization of every
finite Loewy block, with the correct spectral multiplicities.  The issue is
the half-Tate centering.  Put
\begin{equation}
                         S_t=e^{t/2}T_t.                         \tag{6}
\end{equation}
On $k_\rho$, its adjoint character is
$e^{-t(\overline\rho-1/2)}$.

\section{The centering obstruction}

\begin{theorem}[Tate centering is not a finite-rank correction --- \PROVED]
Suppose a positive Hilbert inner product on the full functional pair
$\rho,1-\overline\rho$ makes the centered action (6) contractive and
preserves both spectral characters.  Then
\begin{equation}
                              \Re\rho=\frac12.                   \tag{7}
\end{equation}
Moreover, changing the inner product or adjoining/removing the two polar
Tate lines cannot repair an off-critical pair.
\end{theorem}

\begin{proof}
On the paired semisimple block the centered action has eigenvalues
\[
 e^{-t(\rho-1/2)},\qquad e^{t(\overline\rho-1/2)},
\]
whose determinant has modulus one.  A contraction similar to this block
has both singular values at most one, while their product is one; it is
therefore unitary and both eigenvalues have modulus one.  This is (7).
The argument permits every positive Gram matrix on the pair, so changing
the metric cannot help.  The primitive form removes the two polar Tate
lines before testing the odd pair, so a rank-two polar correction does not
alter this determinant argument.
\end{proof}

\begin{corollary}[Verdict on the completed inner scattering --- \PROVED]
The model space (3) constructs an unconditional positive Hilbert
realization with the correct finite spectral traces, but it does not prove
the centered Witt contractivity.  Promoting $T_t$ to the half-Tate centered
semigroup $S_t$ is contractive on every functional pair if and only if all
those pairs lie on the critical line.
\end{corollary}

Thus the scalar completed scattering solves the existence and trace side of
the Hilbert problem but not row (d).  Its failure is of infinite spectral
rank if there are off-critical pairs; it cannot be absorbed by the two Tate
terminals.  A successful cross-place construction must therefore derive a
new positive conservation law for the centered semigroup, not merely use
the already-inner uncentered scattering (1).

\begin{center}
\begin{tabular}{>{\raggedright\arraybackslash}p{.56\linewidth}
                >{\centering\arraybackslash}p{.15\linewidth}
                >{\raggedright\arraybackslash}p{.19\linewidth}}
\toprule
Statement & Status & Location \\
\midrule
Completed Euler--Gamma inner scattering $\Theta$ & \PROVED & (1)--(3) \\
Positive Hilbert realization of finite zero blocks and traces & \PROVED & (4)--(5) \\
Half-Tate centered contractivity & \OPEN & equivalent to (7) for all zeros \\
Finite Tate correction closes the centered gap & false & Theorem 3 \\
Condition $G$, row (d) & \OPEN & not supplied by uncentered innerness \\
\bottomrule
\end{tabular}
\end{center}

\end{document}
